\chapter{Heisenberg Algebra and the Exact Schwinger Term}
\label{ch:heisenberg_algebra}

\section{Introduction and the Algebraic Origin of the Schwinger Anomaly}

The central mathematical miracle of quantum bosonization is that bilinear products of fermionic field operators satisfy bosonic canonical commutation relations. In continuous quantum field theory, this phenomenon is expressed by the affine Lie algebra $\widehat{\mathfrak{u}}(1)$ (the $U(1)$ Kac--Moody current algebra):
\begin{equation}
[J_m, J_n] = m\, \delta_{m+n, 0}\, I.
\end{equation}
The central term $m\, \delta_{m+n, 0}$, known as the \emph{Schwinger term} or the chiral anomaly, is traditionally derived via heuristic point-splitting, formal normal ordering, or infinite Dirac sea subtractions.

On our discrete lattice $\Lambda = \mathbb{Z}/L\mathbb{Z}$ ($L = 2h$, $h \ge 1$), the situation is fundamentally clarified. 
As proven by the No-Go Theorem~\ref{thm:no_go}, no finite-dimensional matrix algebra can admit operators satisfying $[A, B] = I$ because the trace of any commutator must vanish identically. How, then, can fermionic density operators on the finite Fock space $\mathcal{F}(L)$ satisfy bosonic commutation relations?

\begin{remark}[Energy Budget Regularization]
The answer lies in the \emph{Energy Budget regularization}:
\begin{enumerate}
    \item On the full Fock space $\mathcal{F}(L)$, the commutator $[\rho_{-m}, \rho_m]$ is not a scalar, but a non-trivial dynamical operator given by the difference between the fermion occupations at the bottom and top edges of the band.
    \item When evaluated on vectors in the energy budget subspace $B(N, K)$ satisfying explicit \emph{margin inequalities}, the deep Fermi sea is frozen full and the high momentum band is frozen empty.
    \item Consequently, the edge operators evaluate to the exact positive integer $m$ on budget states, realizing the exact Kac--Moody algebra without any infinite-dimensional or topological approximations.
\end{enumerate}
\end{remark}

\section{The Exact Diagonal Edge Formula on Fock Space}

We first evaluate the commutator of opposing density modes on the entire Fock space $\mathcal{F}(L)$ without any subspace restriction.

Recall that $\rho_m = d\Gamma(T_m)$ is the second quantization of the one-particle partial shift matrix $T_m$ on $\ell^2(\Lambda^*)$ (Chapter~\ref{ch:density_modes}). By Lemma~\ref{lem:dgamma_comm}, second quantization preserves Lie brackets:
\begin{equation}
[\rho_{-m}, \rho_m] = d\Gamma([T_{-m}, T_m]).
\end{equation}

\begin{theorem}[Exact Diagonal Edge Formula]
\label{thm:diagonal_edge_formula}
Let $L = 2h$ with $h \ge 1$. For any integer $1 \le m \le h$, the algebraic commutator $[\rho_{-m}, \rho_m]$ on the entire Fock space $\mathcal{F}(L)$ satisfies the exact operator identity:
\begin{equation}
[\rho_{-m}, \rho_m] = \sum_{q = -h+1}^{-h+m} n_q - \sum_{q = h-m+1}^{h} n_q,
\end{equation}
where $n_q \coloneqq c_q^\dagger c_q$ is the momentum occupation number operator at mode $q$.
\end{theorem}

\begin{proof}
The matrices $T_{-m}T_m$ and $T_mT_{-m}$ are diagonal: the first is the identity on modes $q\le h-m$, while the second is the identity on modes $q\ge-h+1+m$. Their diagonal entries cancel on the bulk interval $-h+1+m\le q\le h-m$; the difference is $+1$ on the bottom edge and $-1$ on the top edge, with all off-diagonal entries zero. Thus
\begin{equation}
[T_{-m},T_m]=\sum_{q=-h+1}^{-h+m}E_{q,q}-\sum_{q=h-m+1}^{h}E_{q,q}.
\end{equation}
Applying $d\Gamma(E_{q,q})=c_q^\dagger c_q=n_q$ gives the asserted identity.
\end{proof}

\begin{definition}[Edge Domains]
For any mode $1 \le m \le h$, the \emph{Bottom Edge Domain} $\mathcal{E}_{\mathrm{bot}}(m)$ and \emph{Top Edge Domain} $\mathcal{E}_{\mathrm{top}}(m)$ are the subsets of $\Lambda^*$:
\begin{align}
\mathcal{E}_{\mathrm{bot}}(m) &\coloneqq \{ q \in \Lambda^* \mid q \le -h + m \} = \{-h+1, -h+2, \dots, -h+m\}, \\
\mathcal{E}_{\mathrm{top}}(m) &\coloneqq \{ q \in \Lambda^* \mid q > h - m \} = \{h - m + 1, \dots, h\}.
\end{align}
Both sets contain exactly $m$ momentum modes: $\#\mathcal{E}_{\mathrm{bot}}(m) = \#\mathcal{E}_{\mathrm{top}}(m) = m$.
\end{definition}

\section{Regime M1: The Diagonal Buffer Zone and the Schwinger Term}

We now evaluate the edge formula on states belonging to the energy budget subspace $B(N, K)$.

\begin{theorem}[Regime M1: Exact Diagonal Schwinger Term]
\label{thm:regime_m1}
Let $\psi \in B(N, K)$ be any state in the energy budget with charge $N \in \{-h, \dots, h\}$ and excitation energy at most $K \ge 0$.
If the mode $m \ge 1$ satisfies the \textbf{M1 Margin Condition}:
\begin{equation}
m + K + |N| \le h,
\end{equation}
then the commutator $[\rho_{-m}, \rho_m]$ evaluates on $\psi$ to the exact scalar Schwinger term:
\begin{equation}
[\rho_{-m}, \rho_m] \psi = m \psi.
\end{equation}
\end{theorem}

\begin{proof}
It suffices by linearity to take a basis vector $|S\rangle$ in the budget. The edge formula gives
\begin{equation}
[\rho_{-m},\rho_m]|S\rangle=\left(\sum_{q\in\mathcal{E}_{\mathrm{bot}}(m)}n_q-\sum_{q\in\mathcal{E}_{\mathrm{top}}(m)}n_q\right)|S\rangle.
\end{equation}
The M1 condition places every bottom edge mode at or below the deep-sea rigidity threshold $N-K$, so all $m$ such modes are occupied; it places every top edge mode strictly above the high-momentum threshold $N+K$, so those $m$ modes are empty. The displayed edge sum is therefore $(m-0)|S\rangle=m|S\rangle$. Linearity gives the result for every $\psi\in B(N,K)$.
\end{proof}

\section{Regime M2: Off-Diagonal Suppression and Pauli Blocking}

We now examine off-diagonal commutators $[\rho_m, \rho_n]$ where $m + n \neq 0$.

\begin{theorem}[Regime M2: Off-Diagonal Suppression]
\label{thm:regime_m2}
Let $\psi \in B(N, K)$. For any mode transfers $m, n \in \mathbb{Z}$ with $m + n \neq 0$, if the parameters satisfy the \textbf{M2 Margin Condition}:
\begin{equation}
|m| + |n| + K + |N| \le h,
\end{equation}
then the density modes commute identically on $\psi$:
\begin{equation}
[\rho_m, \rho_n] \psi = 0.
\end{equation}
\end{theorem}

\begin{proof}
If $m,n$ have the same sign, the global same-sign commutativity theorem applies. For opposite signs, assume without loss of generality $m>0$ and $n=-k<0$, with $m\ne k$. The one-particle commutator consists of hops across the two band edges. Under the M2 margin, each such hop either has its source in the rigidly empty high-momentum region, or its destination in the rigidly occupied deep sea. In the first case annihilation kills the source particle; in the second Pauli exclusion blocks creation at the destination. Thus every edge hop annihilates each budget basis vector, and hence $[\rho_m,\rho_n]\psi=0$ throughout $B(N,K)$.
\end{proof}

\section[The Exact U(1) Kac--Moody Algebra on the Budget Subspace]{The Exact $U(1)$ Kac--Moody Algebra on the Budget Subspace}

Combining Regimes M1 and M2 yields the central theorem of constructive bosonization.

\begin{theorem}[$U(1)$ Kac--Moody Current Algebra on $B(N, K)$]
\label{thm:kac_moody}
Let $M \ge 1$ be a mode cutoff. Under the uniform margin condition:
\begin{equation}
2M + K + |N| \le h,
\end{equation}
the normal-ordered density modes $:\!\rho_m\!:$ satisfy the exact Kac--Moody algebra for all $|m|, |n| \le M$ on the budget subspace $B(N, K)$:
\begin{equation}
\left[ :\!\rho_m\!:, :\!\rho_n\!: \right] \psi = -m\, \delta_{m+n, 0}\, \psi \quad \text{for all } \psi \in B(N, K).
\end{equation}
Equivalently, setting $n = -m$ with $m \ge 1$:
\begin{equation}
\left[ :\!\rho_{-m}\!:, :\!\rho_m\!: \right] \psi = m \psi.
\end{equation}
\end{theorem}

\begin{proof}
If either mode is zero, use $:\!\rho_0\!:=N$ and charge covariance to obtain a zero commutator. When $m+n\ne0$, the uniform margin implies the M2 condition, so Theorem~\ref{thm:regime_m2} gives zero. When $m+n=0$, M1 applies to the positive mode $k=|m|$: if $m<0$, it gives the stated commutator directly, while if $m>0$, antisymmetry reverses the order and gives the stated negative sign. These cases establish the formula.
\end{proof}

\begin{definition}[Harmonic Oscillator Bosonic Modes]
For each mode $m \ge 1$, we define the normalized bosonic creation and annihilation operators on $B(N, K)$ by:
\begin{equation}
b_m \coloneqq \frac{1}{\sqrt{m}} :\!\rho_{-m}\!:, \qquad b_m^\dagger \coloneqq \frac{1}{\sqrt{m}} :\!\rho_m\!: = (b_m)^\dagger.
\end{equation}
Under the margin condition of Theorem~\ref{thm:kac_moody}, these operators satisfy the standard harmonic oscillator algebra on $B(N, K)$:
\begin{equation}
[b_m, b_n^\dagger] \psi = \delta_{m, n} \psi, \qquad [b_m, b_n] \psi = 0, \qquad [b_m^\dagger, b_n^\dagger] \psi = 0.
\end{equation}
\end{definition}

\section{Word-Excursion Tracking in Composed Products}

When evaluating products of density operators $\rho_{a_1} \rho_{a_2} \cdots \rho_{a_s}$ (such as in vertex operators or multi-mode partition states), intermediate states undergo transient energy fluctuations.

\begin{remark}[Interpreting Energy Excursions]
For example, the written word $W=\rho_{-2}\rho_3\rho_{-1}$ acts first by $\rho_{-1}$, then $\rho_3$, then $\rho_{-2}$. The successive energy changes are $-1,3,-2$, with partial sums $-1,2,0$, so the largest upward excursion is $2$. This maximum tells us how much extra energy budget must be available for every intermediate state, even though the final net change is zero.
\end{remark}

\begin{definition}[Application Order and Cumulative Energy Excursion]
For a written operator product $W = \rho_{a_1} \rho_{a_2} \cdots \rho_{a_s}$, the operators act from right to left in the sequence of signed transfers $(a_s, a_{s-1}, \dots, a_1)$.
\begin{enumerate}
    \item The partial energy shift after applying the first $r$ operators ($0 \le r \le s$) is:
    \begin{equation}
    \Delta E_r \coloneqq \sum_{j=1}^r a_{s - j + 1}.
    \end{equation}
    \item The \emph{maximum upward energy excursion} $X(W)$ is:
    \begin{equation}
    X(W) \coloneqq \max\left(0, \; \max_{0 \le r \le s} \Delta E_r \right).
    \end{equation}
\end{enumerate}
\end{definition}

\begin{theorem}[Valid Margin Excursion Accounting]
\label{thm:excursion_accounting}
Let $W = \rho_{a_1} \dots \rho_{a_s}$ be an operator word with mode cutoff $|a_j| \le M$. If the input state $\psi \in B(N, K)$ satisfies:
\begin{equation}
2M + K + X(W) + |N| \le h,
\end{equation}
then every intermediate state generated during the evaluation of $W \psi$ belongs to $B(N, K + X(W))$, and every commutator $[:\!\rho_m\!:, :\!\rho_n\!:]$ encountered along the evaluation acts as the exact scalar $-m \delta_{m+n, 0}$.
\end{theorem}

\begin{proof}
By Theorem~\ref{thm:budget_action_annihilation}, applying $\rho_{a_s}$ to $\psi \in B(N, K)$ produces a state in $B(N, K + a_s)$. By induction, after $r$ steps the state belongs to $B(N, K + \Delta E_r) \subseteq B(N, K + X(W))$.
Since $2M + (K + X(W)) + |N| \le h$, each intermediate state satisfies the uniform margin condition of Theorem~\ref{thm:kac_moody}. Therefore, all intermediate commutators evaluate to exact scalars.
\end{proof}

\begin{remark}[Budget-Local Resolution of the Trace Obstruction]
On the full finite Fock space, $[\rho_{-m},\rho_m]$ is the difference of the bottom- and top-edge occupation sums, so its trace vanishes and it is not the scalar $mI$. On a budget subspace with the M1 margins, the bottom edge is fully occupied and the top edge empty; the same operator therefore evaluates as $mI$ there. The central term is thus a local identity on the physically selected low-energy states, consistent with the finite-dimensional trace constraint.
\end{remark}

\section{Chapter Summary}
The central result is a budget-local Heisenberg relation for the finite
density modes. It resolves the apparent conflict with the trace identity:
the commutator is traceless on the full finite Fock space, but evaluates to
the central scalar on the low-energy states protected by the margin
condition. Tracking intermediate energy excursions explains exactly when
products and reordered words remain inside that valid region. This gives a
finite, domain-conscious counterpart to the continuum current algebra and
sets up the partition-counting proof of Haldane completeness.
